
\subsubsection{Interpretation of Findings}

Variance measurements establish that seed-based variance in simple MLPs is measurable, task-dependent, and mechanistically explained through a validated causal chain. The $N\geq 30$ criterion from Rajput and Kumar (2023) holds for variance detection (statistical significance testing) but requires refinement for precision (stable quantitative estimation).

Task difficulty influences variance magnitude. The 9-$10\times$ scaling between MNIST (0.04\%) and Fashion-MNIST (0.35-0.59\%) reflects a ceiling constraint. MNIST achieves >97\% accuracy, compressing variance into <3\% absolute range. Fashion-MNIST achieves ~88\% accuracy, allowing larger absolute variance. This suggests that variance measurement is practical for tasks with 80-95\% baseline accuracy but ceiling-compressed for tasks above 95\%.

The complete causal chain (seed $\rightarrow$ independent weights $\rightarrow$ different trajectories $\rightarrow$ different minima $\rightarrow$ measurable variance) is supported by statistical evidence (p < $10^{-6}$ at each step). This mechanistic understanding enables predictions: deeper architectures will show higher variance (confirmed: 2-layer shows 1.5-1.$7\times$ increase), and tasks near accuracy ceiling will show compressed variance (confirmed: MNIST 0.04\% vs. Fashion-MNIST 0.35-0.59\%).

The detection-vs-precision boundary indicates that N=30 provides sufficient power for detecting non-zero variance (p<0.05) but insufficient stability for precise quantification (bootstrap CI widths 93-110\% vs. 50\% threshold). This finding refines existing theory: sample size requirements differ for hypothesis testing (N=30) vs. narrow estimation (likely N>50 for neural networks).

\subsubsection{Limitations}

\textbf{Limitation 1: N=30 insufficient for bootstrap precision}

Bootstrap confidence interval widths (93-110\%) exceed the 50\% stability threshold. Users of this baseline must acknowledge measurement uncertainty when citing specific variance values or increase sample size for narrower estimates. The limitation is quantitative (collect more seeds) rather than qualitative (variance is unmeasurable). Variance detection remains validated (p<0.05).

Future work could conduct N sensitivity analysis by replicating experiments with N $\in$ {50, 100, 200} and plotting CI width vs. N to identify the threshold where width $\leq 50$\%.

\textbf{Limitation 2: MNIST ceiling effect}

MNIST achieves 98\% accuracy, leaving minimal absolute range (<2\%) for cross-seed variance. Measured variance (0.04\%) falls below the practical detectability threshold (0.3\%). The hypothesis is confirmed for medium-difficulty tasks but not applicable to easy tasks above 95\% accuracy.

The dual-dataset design (MNIST + Fashion-MNIST) was intended to test task-dependency. Discovering 9-$10\times$ scaling between easy and medium tasks is a finding rather than a failure. Future work could extend to intermediate-difficulty datasets (KMNIST ~92\%, EMNIST ~89\%) to map variance vs. accuracy relationships systematically.

\textbf{Limitation 3: Architecture scope limited to simple MLPs}

Validation is limited to 1-layer (101,770 parameters) and 2-layer (235,146 parameters) fully-connected networks. Variance dynamics may differ for deeper architectures (CNNs, ResNets, Transformers). The observed 1.5-1.$7\times$ variance increase (2-layer vs. 1-layer) may not extrapolate linearly to 50-layer ResNets.

The scope decision establishes the simplest baseline before scaling to complex architectures. Picard et al. (2021) demonstrated feasibility on CIFAR-10 ResNet-18, providing an existence proof for CNN extension. Future work could validate the protocol on CIFAR-10 CNNs.

\textbf{Limitation 4: Bootstrap i.i.d. assumption potentially violated}

Bootstrap resampling assumes independent and identically distributed samples. The 30 training runs share dataset, optimizer, and architecture, with only random seed differing. This may violate the i.i.d. assumption, potentially contributing to the 93-110\% CI widths.

The limitation affects only precision measurement (H-M3), not core variance detection findings (H-E1, H-M1, H-M2). Future work could compare bootstrap CIs with permutation test confidence intervals and Bayesian hierarchical model credible intervals on existing data to assess whether i.i.d. violation contributes to wide CIs.

\subsubsection{Implications for Uncertainty Quantification Research}

This work provides calibration infrastructure for uncertainty quantification methods. Researchers developing MC Dropout, Bayesian NNs, or ensemble methods can measure seed-based variance $\sigma ^{2}_seed$ using the validated protocol and compare against method-specific variance $\sigma ^{2}_method$. If $\sigma ^{2}_method \approx \sigma ^{2}_seed$, the method captures only initialization uncertainty. If $\sigma ^{2}_method$ > $\sigma ^{2}_seed$, the method adds epistemic/aleatoric uncertainty beyond initialization.

\subsubsection{Implications for Reproducibility}

Practitioners reporting mean $\pm$ std across N seeds can use the following guidelines:
\begin{itemize}
\item N=30 for significance testing (p<0.05) to confirm variance exists
\item N>50 for precise quantification (bootstrap CI $\leq 50$\%) based on extrapolation
\item N=3-5 (common practice) is likely underpowered for both detection and precision
\end{itemize}

Variance benchmarks by task difficulty:
\begin{itemize}
\item Easy tasks (>95\% accuracy): Expect $\sigma \leq$ 0.1\% (ceiling effect)
\item Medium tasks (80-95\% accuracy): Expect $\sigma$ = 0.3-0.8\%
\item Hard tasks (<80\% accuracy): Unknown (beyond validated range)
\end{itemize}

\subsubsection{Future Work}

Immediate extensions include:
\begin{enumerate}
\item N sensitivity analysis to validate the detection-vs-precision boundary by testing N $\in$ {50, 100, 200}
\item Statistical triangulation comparing bootstrap, permutation test, and Bayesian methods to validate i.i.d. assumption robustness
\item Task difficulty gradient experiments adding intermediate datasets (KMNIST, EMNIST) to map variance vs. accuracy relationships
\end{enumerate}

Longer-term directions include:
\begin{itemize}
\item Comprehensive variance atlas across architectures (MLPs, CNNs, ResNets, Transformers) and tasks
\item MC Dropout calibration study demonstrating baseline usage
\item Integration into MLOps pipelines for deployment reliability monitoring
\end{itemize}
